\section{From $W(\mathsf{SC}^{\mathrm{ch,top}})$ to axioms and back}
\label{sec:operad-to-axioms-equiv}
\label{sec:equivalence-chapter}

The boundary chart $A_b = \End_\cC(b)$ of the open-sector factorization category admits two equivalent descriptions: operadic ($W(\SCchtop)$-algebra structure) and axiomatic (the explicit $\Ainf$ chiral operations $\{m_k\}$).  The operadic formulation is canonical; the axiomatic formulation is computable.

\subsection{From the operad to explicit $m_k$ with spectral substitution}
\label{subsec:operad-to-mk-equiv}
Let $(A_{\mathsf{ch}},A_{\mathsf{top}})$ be a $C_\ast\!\bigl(W(\mathsf{SC}^{\mathrm{ch,top}})\bigr)$–algebra with underlying complex $A$ and translation operator $\partial$ (from the holomorphic \(\C\)–action). Choose, for each $k\ge1$, a fundamental chain $[\FM_k(\C)]\in C_\ast(\FM_k(\C))$ and similarly for $E_1$ (recording the $t$–ordering). Define
\[
m_k(a_1,\dots,a_k)
:= \bigl\langle\ \rho\bigl([\FM_k(\C)]\times [E_1(\text{ordering})]\bigr)\,,\, a_1\otimes\cdots\otimes a_k\ \bigr\rangle
\ \in\ A((\lambda_1))\cdots((\lambda_{k-1})),
\]
where $\rho$ is the operad action and the \(\lambda\)–dependence is the generating function of the log‑forms (residue degrees) on $\FM_k(\C)$ prescribed by the operad.
Then:
\begin{enumerate}[label=(\alph*)]
  \item sesquilinearity \eqref{eq:sesqui-left}–\eqref{eq:sesqui-right} follows from equivariance under translations in $\C$ (the \(\C\)–action generates $\partial$ and gives \(\lambda\)-multiplications),
  \item the spectral substitution law in \eqref{eq:ainfty-relation-raw} is the image of the boundary identification of faces in $\partial\FM_k(\C)$ where a consecutive block collapses (the effective parameter is the sum of inner ones).
\end{enumerate}

\subsection{Rectification}
\label{subsec:rectification-equiv}
Homotopy-Koszulity of $\mathsf{SC}^{\mathrm{ch,top}}$
(Theorem~\ref{thm:homotopy-Koszul}) upgrades the bar-cobar
adjunction to a Quillen equivalence
(Theorem~\ref{thm:bar-cobar-adjunction}), guaranteeing
essential uniqueness.  The full four-step proof appears in
Theorem~\ref{thm:rectification}.

\subsection{Compatibility with the PVA construction}
\label{subsec:compat-PVA-equiv}
The proof of the PVA structure on cohomology in Section~\ref{sec:axioms-Ainfty-chiral} can be interpreted operadically as the statement that the associated graded (holomorphic filtration) of $W(\mathsf{SC}^{\mathrm{ch,top}})$ is the free product of BD chiral and $E_1$, so the cohomological operations descend from the closed color; the bracket is the classical limit (residue) of the operadic composites, and associativity is the regular part of the $m_2$ identity.

